% Stage 8 restored fragment — 2.5.4 Radiative and non-radiative components (archive lines 7127-7360)
\subsection{Radiative and non-radiative components}
\label{sec:ch254}
A compact-range engineer can measure strong electric and magnetic fields within one wavelength
of a horn mouth while the far-zone gain remains modest: reactive energy stored near the feed
is large, yet little of it crosses a distant measurement sphere as export
\cite{poynting1884}. Conversely, currents in $\ker\mathcal{R}_\omega$ can sustain
elaborate near-field structure while producing no far-field signature---the classical
non-radiating source \cite{harrington2001}. Inverse synthesis and scattering campaigns inherit
this split: near-field scans and angle-resolved data constrain only the radiative image
$\mathcal{R}_\omega[\Jimp]$, leaving non-radiative content invisible until additional
excitation or realizability constraints are imposed; Chapter~\ref{ch:04} develops that
$\mathcal{S}_{\mathrm{real}}$ question spectrally \cite{stratton1941}.

The component partition below imports reactive--radiative logic from
Section~\ref{sec:ch01.2.4} and spectral projectors
Eqs.~\eqref{eq:ch2.spectral-projectors}--\eqref{eq:ch2.rad-null-on-evanescent} without
re-proving projector algebra \cite{stratton1941}. Complete non-radiating null-space
classification belongs to Section~\ref{sec:ch27}.
Let \(\mathbf{u}_\omega=(\E,\H)=\mathcal{S}_\omega[\mathfrak{D}]\) and
\(\mathbf{u}_{\mathrm{rad}}=\mathcal{R}_\omega[\mathfrak{D}]\) be as in
Eq.~\eqref{eq:ch2.section25-evaluation-chain} \cite{stratton1941}. Spectral
projectors \(\Pi_{\mathrm{prop}}\), \(\Pi_{\mathrm{ev}}\), \(\Pi_{\mathrm{non}}\) act on
\(\widehat{\mathcal{F}}_\Omega\) as in Subsection~\ref{sec:ch224}; outgoing selector
\(\mathcal{O}_{\mathrm{out}}\) and radiation map satisfy
Eqs.~\eqref{eq:ch2.outgoing-radiation-refinement}--\eqref{eq:ch2.rad-projection-composition}
\cite{stratton1941}. Regime qualification uses
\(\mathcal{D}_{J,\mathrm{reg}}\) from Eq.~\eqref{eq:ch1.regime-qualified-decision} and
criterion Eq.~\eqref{eq:ch2.regime-operator-screen} \cite{stratton1941}.
On admissible \(\mathfrak{D}\),
\begin{equation}
\label{eq:ch2.integral-spectral-decomposition}
\mathcal{S}_\omega[\mathfrak{D}]
=
\Pi_{\mathrm{prop}}\mathcal{S}_\omega[\mathfrak{D}]
+\Pi_{\mathrm{ev}}\mathcal{S}_\omega[\mathfrak{D}]
+\Pi_{\mathrm{non}}\mathcal{S}_\omega[\mathfrak{D}],
\end{equation}
restating Eq.~\eqref{eq:ch2.propagating-evanescent-decomposition} on the integral outputs
constructed above \cite{stratton1941}.
Define the \emph{non-propagating} component
\begin{equation}
\label{eq:ch2.nonpropagating-component-def}
\mathbf{u}_{\mathrm{np}}
=
\Pi_{\mathrm{ev}}\mathcal{S}_\omega[\mathfrak{D}]
+\Pi_{\mathrm{non}}\mathcal{S}_\omega[\mathfrak{D}],
\end{equation}
and the \emph{propagating} component
\(\mathbf{u}_{\mathrm{prop}}=\Pi_{\mathrm{prop}}\mathcal{S}_\omega[\mathfrak{D}]\)
\cite{stratton1941}. Only \(\mathbf{u}_{\mathrm{prop}}\) is eligible for outgoing radiation
selection; evanescent and non-propagating content satisfies
Eq.~\eqref{eq:ch2.rad-null-on-evanescent} \cite{harrington2001}. Equation
\eqref{eq:ch2.integral-spectral-decomposition} applies the \ref{ch:01} partition to
integral outputs: reactive dominance on near shells corresponds to large
\(\|\mathbf{u}_{\mathrm{np}}\|\) relative to \(\|\mathbf{u}_{\mathrm{prop}}\|\) in the sense
of Eq.~\eqref{eq:ch2.subspace-energy-inequality} \cite{stratton1941}.
The radiative component is
\begin{equation}
\label{eq:ch2.radiative-component-def}
\mathbf{u}_{\mathrm{rad}}
=
\mathcal{R}_\omega[\mathfrak{D}]
=
\Pi_{\mathrm{rad}}\,\mathcal{S}_\omega[\mathfrak{D}]
=
\mathcal{O}_{\mathrm{out}}\circ\Pi_{\mathrm{prop}}\,\mathcal{S}_\omega[\mathfrak{D}],
\end{equation}
using Eqs.~\eqref{eq:ch2.outgoing-radiation-refinement} and~\eqref{eq:ch2.rad-projection-composition}
\cite{stratton1941}. The complementary \emph{non-radiative} component is
\begin{equation}
\label{eq:ch2.nonradiative-component-def}
\mathbf{u}_{\mathrm{nr}}
=
\mathcal{S}_\omega[\mathfrak{D}]-\mathbf{u}_{\mathrm{rad}}
=
\big(\mathrm{id}-\Pi_{\mathrm{rad}}\big)\mathcal{S}_\omega[\mathfrak{D}].
\end{equation}
Equations~\eqref{eq:ch2.radiative-component-def}--\eqref{eq:ch2.nonradiative-component-def}
partition every admissible field pair into export and non-export parts without altering
Maxwell dynamics: \(\mathbf{u}_\omega=\mathbf{u}_{\mathrm{rad}}+\mathbf{u}_{\mathrm{nr}}\)
\cite{stratton1941}. Non-radiative content includes evanescent, non-propagating, and
incoming-propagating fractions annihilated by $\Pi_{\mathrm{rad}}$; large near-field
magnitude therefore does not establish export \cite{stratton1941}. Sources in
$\ker\mathcal{R}_\omega$ produce $\mathbf{u}_{\mathrm{rad}}=\mathbf{0}$ while
$\mathbf{u}_{\mathrm{nr}}\neq\mathbf{0}$ (Eq.~\eqref{eq:ch2.ker-R-omega})
\cite{stratton1941}.
\begin{figure}[t]
\centering
\ChOneFigBlock{%
\begin{tikzpicture}[ch1fig, node distance=5mm and 7mm]
 \node[ch1box] (s) {$\mathcal{S}_\omega[\mathfrak{D}]$};
 \node[ch1box, right=7mm of s] (pp) {$\Pi_{\mathrm{prop}}$};
 \node[ch1box, right=7mm of pp] (pr) {$\Pi_{\mathrm{rad}}$};
 \node[ch1box, below=5mm of s] (np) {$\mathbf{u}_{\mathrm{nr}}$};
 \node[ch1box, below=5mm of pp] (ev) {$\Pi_{\mathrm{ev}}+\Pi_{\mathrm{non}}$};
 \draw[ch1flow] (s) -- (pp) -- (pr);
 \draw[ch1flow] (s.south) -- (s.south |- np.north);
 \draw[ch1flow] (s.south |- np.north) -- (np.north);
 \draw[ch1flow] (s.east) -- (s.east |- ev.west);
 \draw[ch1flow] (s.east |- ev.west) -- (ev.west);
\end{tikzpicture}%
}
\caption{Radiative/non-radiative split: \(\mathbf{u}_{\mathrm{rad}}\) is
\(\Pi_{\mathrm{rad}}\mathcal{S}_\omega\); \(\mathbf{u}_{\mathrm{nr}}\) collects all
non-export content including Eq.~\eqref{eq:ch2.nonpropagating-component-def}.}
\label{fig:ch2.rad-nr-component-split}
\end{figure}
Figure~\ref{fig:ch2.spectral-radiation-criterion} from Subsection~\ref{sec:ch224} on integral
outputs \cite{stratton1941}. Using energy norm Eq.~\eqref{eq:ch2.field-energy-norm}, define
\begin{equation}
\label{eq:ch2.radiative-energy-functional}
\mathcal{E}_{\mathrm{rad}}[\mathfrak{D}]
=
\|\mathbf{u}_{\mathrm{rad}}\|_{\mathcal{E}}^{2},
\qquad
\mathcal{E}_{\mathrm{nr}}[\mathfrak{D}]
=
\|\mathbf{u}_{\mathrm{nr}}\|_{\mathcal{E}}^{2},
\end{equation}
with \(\|\cdot\|_{\mathcal{E}}\) from Eq.~\eqref{eq:ch2.radiation-energy-norm} on each
component \cite{stratton1941}. Shell power splits analogously:
\begin{equation}
\label{eq:ch2.component-power-budget}
\mathcal{P}_{\mathrm{tot}}(r)
=
\mathcal{P}_{\mathrm{rad}}(r;\Jimp)
+\mathcal{P}_{\mathrm{nr}}(r;\Jimp),
\end{equation}
where \(\mathcal{P}_{\mathrm{rad}}\) is Eq.~\eqref{eq:ch2.power-flux-functional} and
\begin{equation}
\label{eq:ch2.nonradiative-power-functional}
\mathcal{P}_{\mathrm{nr}}(r;\Jimp)
=
\Re\oint_{S_r}\hat{\mathbf{r}}\cdot\mathbf{S}_{\mathrm{nr}}[\Jimp]\,dS,
\qquad
\mathbf{S}_{\mathrm{nr}}
=
\tfrac12\,\E_{\mathrm{nr}}\times\H_{\mathrm{nr}}^{*},
\end{equation}
with \((\E_{\mathrm{nr}},\H_{\mathrm{nr}})=\mathbf{u}_{\mathrm{nr}}\)
\cite{stratton1941}. Equation~\eqref{eq:ch2.component-power-budget} is not
a new conservation law; it decomposes Eq.~\eqref{eq:ch2.radiated-power-screen} when both
components are evaluated on the same shell \cite{stratton1941}. In the far zone,
\(\mathcal{P}_{\mathrm{nr}}(r)\to 0\) and \(\mathcal{P}_{\mathrm{rad}}(r)\to P_\infty\) under
Eq.~\eqref{eq:ch1.radiative-regime-condition} \cite{stratton1941}. Near-zone
shells with \(\chi(r)\gg 1\) from Eq.~\eqref{eq:ch1.regime-index} typically have
\(|\mathcal{P}_{\mathrm{nr}}|\gg|\mathcal{P}_{\mathrm{rad}}|\) \cite{stratton1941}.
Regime qualification from Section~\ref{sec:ch01.2.4} becomes
\begin{equation}
\label{eq:ch2.regime-component-criterion}
\mathcal{D}_{J,\mathrm{reg}}=1
\ \Longrightarrow\
\text{transport uses}\ \mathbf{u}_{\mathrm{rad}},\
\mathcal{P}_{\mathrm{rad}},\
\mathbf{M}_{\mathrm{rad}}[\Jimp],
\end{equation}
and forbids interpreting \(\mathbf{u}_{\mathrm{nr}}\), \(\mathcal{P}_{\mathrm{nr}}\), or
\(\Pi_{\mathrm{ev}}\mathcal{S}_\omega[\mathfrak{D}]\) as export channels
\cite{stratton1941}. Equation~\eqref{eq:ch2.regime-component-criterion} is the
component form of Eq.~\eqref{eq:ch2.regime-operator-screen} after
Eqs.~\eqref{eq:ch2.radiative-component-def}--\eqref{eq:ch2.nonradiative-component-def}.
Laboratory errors from Section~\ref{sec:ch01.2.4}---near-zone reactive scans, cutoff-guide
modes, evanescent dominance---appear as large \(\mathcal{E}_{\mathrm{nr}}\) or
\(\|\Pi_{\mathrm{ev}}\mathcal{S}_\omega[\mathfrak{D}]\|\) with small
\(\mathcal{P}_{\mathrm{rad}}\) \cite{stratton1941}. Link to radiation null space. Non-radiative components encode source-side invisible content:
\begin{equation}
\label{eq:ch2.nr-image-of-null-space}
\Jimp\in\ker\mathcal{R}_\omega
\ \Longrightarrow\
\mathbf{u}_{\mathrm{rad}}=\mathbf{0},
\quad
\mathbf{u}_{\mathrm{nr}}=\mathcal{S}_\omega[\Jimp]\neq\mathbf{0}\ \text{in general},
\end{equation}
as recorded in Eq.~\eqref{eq:ch2.null-radiation} \cite{harrington2001}.
Conversely, \(\mathbf{u}_{\mathrm{nr}}=\mathbf{0}\) does not imply trivial
\(\mathcal{S}_\omega[\mathfrak{D}]\): a pure outgoing propagating state may have
\(\mathbf{u}_{\mathrm{rad}}=\mathbf{u}_\omega\) \cite{stratton1941}. Section~\ref{sec:ch272}
classifies \(\ker\mathcal{R}_\omega\); here only the field-side split
Eq.~\eqref{eq:ch2.nonradiative-component-def} is identified as the operator image of
non-export content \cite{harrington2001}.
\phantomsection\label{prop:ch2.far-zone-radiative-dominance}
Let \(\mathcal{D}_{J,\mathrm{reg}}=1\) and let \(r\) lie in the far zone of
Eq.~\eqref{eq:ch2.distance-regime-definitions}. Then
\begin{equation}
\label{eq:ch2.far-zone-component-dominance}
\frac{\mathcal{E}_{\mathrm{nr}}[\mathfrak{D}]}{\mathcal{E}_{\mathrm{rad}}[\mathfrak{D}]+\epsilon}
\to 0,
\qquad
\frac{|\mathcal{P}_{\mathrm{nr}}(r)|}{\mathcal{P}_{\mathrm{rad}}(r)+\epsilon}\to 0,
\end{equation}
as \(r\to\infty\) for fixed \(\mathfrak{D}\), with \(\epsilon>0\) regularizing zero export
\cite{stratton1941}. Equivalently, \(\mathbf{u}_\omega\to\mathbf{u}_{\mathrm{rad}}\) in
far-zone observables metered by Eqs.~\eqref{eq:ch2.far-field-integral-reduction}--\eqref{eq:ch2.power-flux-functional}.
Eq.~\eqref{eq:ch2.rad-null-on-evanescent} annihilates evanescent and non-propagating
contributions in \(\Pi_{\mathrm{rad}}\). Far-zone asymptotics and flux stabilization on
large shells give \(\mathcal{P}_{\mathrm{rad}}(r)\to P_\infty>0\) while reactive terms decay
relatively \cite{stratton1941}. Regime condition
Eq.~\eqref{eq:ch1.radiative-regime-condition} removes reactive dominance in the limit.
Below guide cutoff Eq.~\eqref{eq:ch2.cutoff-condition}, modal content lies in
\(\Pi_{\mathrm{non}}\mathcal{S}_\omega[\mathfrak{D}]\) or
\(\Pi_{\mathrm{ev}}\mathcal{S}_\omega[\mathfrak{D}]\) \cite{stratton1941}. Applying
Eq.~\eqref{eq:ch2.far-field-amplitude-integral} or transport indices to such modes violates
Eq.~\eqref{eq:ch2.regime-component-criterion} as in the TE\(_{10}\) illustration of
Section~\ref{sec:ch01.1.2} \cite{harrington2001}. Spectral projectors and regime logic must
therefore be applied together on the same integral outputs \cite{stratton1941}.
Equations~\eqref{eq:ch2.integral-spectral-decomposition}--\eqref{eq:ch2.far-zone-component-dominance}
require linear media, fixed \(\omega\), outgoing radiation condition, and spectral completion sufficient
for Eqs.~\eqref{eq:ch2.spectral-projectors}--\eqref{eq:ch2.spectral-projector-algebra}
\cite{stratton1941}. Sharp boundaries between subspaces may blur when
loss, anisotropy, or coupled interfaces mix propagating and evanescent sheets
\cite{stratton1941}. Component splits are diagnostic on shells, not substitutes for
exterior radiative uniqueness (Subsection~\ref{sec:ch243}) on the full field
\cite{harrington2001}.
Outgoing source-to-field integrals map impressed data to evaluable fields; distance-regime
labels fix where asymptotic reductions apply; flux and operator-norm meters act on
$\mathcal{R}_\omega[\Jimp]$ after regime qualification; and the radiative--non-radiative split
Eqs.~\eqref{eq:ch2.radiative-component-def}--\eqref{eq:ch2.regime-component-criterion}
separates export from storage on the same integral outputs \cite{stratton1941}. Section~\ref{sec:ch26}
develops spectral decomposition of radiation operators on the representations constructed here
\cite{harrington2001}.
The radiative--non-radiative split is what separates a wattmeter that has been qualified on $\Gamma_R$ from a near-field probe that still sees storage and evanescent overlap in $\mathcal{S}_\omega[\Jimp]$ \cite{harrington2001}. Both components coexist in one Maxwell solution until projectors and outgoing selection name which part export meters may integrate \cite{poynting1884}. Evanescent sheets must not be dismissed as mere mathematical artifacts; however, Maxwell theory shows they carry reactive energy measurable on enclosing scans while remaining null under far-zone angular accessibility maps developed in Section~\ref{sec:ch29} \cite{stratton1941}. Section~\ref{sec:ch26} now organizes the radiative image itself into eigenmodes and recoverable rank on the integral representations constructed in Sections~\ref{sec:ch23}--\ref{sec:ch25}.
